\chapter{Research Problem and Methodology}
\label{ch:methodology}

\section{Task Formulation}

Let a dataset example be a tuple $(x, S)$ where $x$ is a UTF-8 text string and $S = \{(s_i, \ell_i, b_i, e_i)\}_{i=1}^{n}$ is a set of gold skill spans.
Each span has surface form $s_i$, label $\ell_i \in \{\textsc{Hard\_Skill}, \textsc{Soft\_Skill}\}$, and character offsets $(b_i, e_i)$ such that $x[b_i:e_i] = s_i$.

The generative extractor produces a structured output $Y = \{(t_j, \hat{\ell}_j)\}$ via instruction-conditioned decoding.
Predicted spans are aligned back to $x$ by exact or case-insensitive substring search; unaligned predictions are counted but excluded from span-level true positives.

\section{Generative Training Objective}

Fine-tuning uses causal language modeling on formatted prompt--target pairs.
For example $x$ with gold skills serialized as JSON, the model minimizes next-token negative log-likelihood over the target tokens while the prompt tokens may be masked according to the Hugging Face \texttt{Trainer} configuration in \texttt{train\_sft\_v2.py}.
This is standard supervised fine-tuning (SFT), not a token-level CRF or span classifier head.

\section{Evaluation Metrics}

Following the benchmark evaluation module, the thesis reports:

\begin{itemize}
  \item \textbf{Exact span F1}: span boundary and label must match gold exactly after alignment.
  \item \textbf{Partial span F1}: overlap-based span match with correct label.
  \item \textbf{Token-level F1}: token membership in skill spans.
  \item \textbf{Per-label F1}: separate scores for hard and soft skills.
  \item \textbf{Parse failure rate}: fraction of examples whose raw JSON could not be parsed.
  \item \textbf{Unaligned prediction rate}: fraction of predicted span texts not found in source text.
\end{itemize}

For wrapper variants, \textbf{taxonomy linking coverage} is reported as the fraction of extracted spans assigned a canonical skill identifier by the linker when taxonomy files are available.

\section{Experimental Methodology}

Two experimental phases are distinguished:

\begin{enumerate}
  \item \textbf{Legacy multi-LLM benchmark} (\texttt{SkillLLM\_v2}, commit \texttt{0627abf}): five models, three training epochs each, metrics in \texttt{results/skillspan\_baselines/}.
  \item \textbf{Stage-2 system evaluation} (\texttt{SkillLLM\_v2\_stage2}): Mistral-focused 80-epoch training, encoder baselines, and wrapper ablations, metrics in \texttt{artifacts/skillllm\_v2/}.
\end{enumerate}

Results from these phases are not merged into a single training run table without explicit labeling because epoch counts and LoRA ranks differ between phases.

\section{Non-Regression Wrapper Methodology}

Stage-2 wrapper variants must not degrade the Mistral LoRA anchor beyond configured floors on exact span F1 and soft-skill F1 (\texttt{configs/skillllm\_v2\_wrapper\_guardrails.yaml}).
If a candidate variant fails guardrails, the gate falls back to anchor predictions for extraction metrics while still recording linking outputs where applicable.
